% !TeX root = ../教材题库_全部条目_转写工作稿.tex
% Source: XJ-B1-C04-TH-P132-Y0715
% 第4章 / 4.4 函数与方程 / 教材 p.132 / PDF 第139页
% OCR未校对：公式与文字均可能有误；未完成数学TeX转写；禁止直接发给学生
% 原题及相邻原解截图（核对用）：../题目原图/XJ-B1-C04-TH-P132-Y0715-01.webp
% 本文件只在转写完成后作为题干使用；OCR 条目尚不是可靠题干。
\begin{bankitem}{p.132 思考（含紧邻说明）}{第4章 / 4.4 函数与方程 / 教材 p.132 / PDF 第139页}{XJ-B1-C04-TH-P132-Y0715}{OCR未校对：公式与文字均可能有误；未完成数学TeX转写；禁止直接发给学生}
\begin{ocrdraft}
[教材p.132 / PDF第139页]
\par
？思考
\par
假设从某水库闸房到防洪指挥部的电话线路发生了故障．这是一条
\par
10km长的线路，在这条线路上有200多根电线杆．想一想：维修工人应怎样最合
\par
理地迅速查出故障所在地呢？
\par
如图4．4一4，工人首先从线路的中点C查起，如果CB段正常，就选择CA的
\par
中点D测试；如果DA段正常，就选择DC的中点E继续测试，一“像检修线路所用
\par
的这种方法称作二分法
\par
闸房
\par
待查
\par
指挥部
\par
DC
\par
图4．4一4
\par
二分法还可以用来寻找函数的零点，迅速地缩小搜索范围，接近零点的准确
\par
位置．
\end{ocrdraft}
\end{bankitem}
